% TEMPLATE-MILC-v3.tex - plantilla maestra UNICA de las guias MILC v3.
%
% La usan las cuatro familias de guias, cada una con su driver:
%   · Grados 6-11          -> scripts/build-guias-g11.py
%   · Bebras               -> scripts/build-guias-bebras.py
%   · Semillero            -> scripts/build-guias-semillero.py
%   · Territorio interior  -> scripts/build-guias-territorio-interior.py
%
% Antes habia tres copias casi identicas (~400 lineas iguales, 6 distintas):
% cada mejora habia que aplicarla tres veces, y en la practica no se hacia
% (el motor de recursos vivia solo en dos de ellas). Lo que varia entre
% familias esta parametrizado en 6 placeholders que cada driver llena
% (se nombran aqui SIN su sintaxis real: el builder reemplaza por texto plano,
% tambien dentro de los comentarios, y un valor multilinea rompe el preambulo):
%
%   HEADER_IZQ        encabezado izquierdo de pagina
%   HEADER_DER        encabezado derecho de pagina
%   PORTADA_ETIQUETA  rotulo sobre el numero grande (GRADO/MOMENTO/MODULO)
%   PORTADA_SERIE     linea de serie de la portada
%   PORTADA_ID        identificador de la guia en la portada
%   PORTADA_CIRCULO   el circulito de grado (°), o vacio si no aplica
%   PDF_TITULO        titulo en texto plano (sin \textbf ni comillas TeX) para
%                     los metadatos del PDF (/Title); lo produce lib_xelatex.texto_pdf
%
% Composicion (auditoria 2026-09-03): las bandas de estacion piden espacio
% (\Needspace*) y prohiben el salto justo despues (\nopagebreak), asi no quedan
% huerfanas al pie; las cajas largas (softbox, drawbox, infoband, puentebox) son
% `breakable`, asi una caja de media pagina ya no arrastra todo a la siguiente
% dejando la anterior al 40 %. Todo texto sobre mostaza o verde va en milcNegro
% (blanco sobre #E5B400 da 1,93:1 y sobre #58A923 2,95:1; ninguno pasa WCAG AA).
% El resto de claves las documenta content/guias/_SCHEMA.md. El script valida
% que no queden placeholders sin reemplazar antes de escribir el archivo final.

% Etiquetado PDF (latex-lab): /Lang, estructura y texto alternativo de las
% figuras, para lectores de pantalla. Debe ir ANTES de \documentclass.
\DocumentMetadata{lang=es-CO,tagging=on}
\documentclass[11pt,letterpaper]{article}
% headheight=24pt: la cabecera derecha lleva dos lineas (identidad de la guia
% y estacion actual). headsep se acorta en la misma medida para que el bloque
% de texto quede exactamente donde estaba con headheight=12pt.
\usepackage[margin=2.0cm,headheight=24pt,headsep=13pt]{geometry}
\usepackage[table]{xcolor}
\usepackage{fontspec}
\usepackage[spanish]{babel}
\usepackage{tikz}
% Librerías TikZ para los diagramas de las guías (bloque `recursos:`):
%  · babel  → babel en español vuelve activos caracteres como «>», y TikZ los usa
%             en sus opciones (>=stealth, ->). Sin esta librería, cualquier
%             diagrama con flechas revienta con «\language@active@arg> has an
%             extra }». Debe ir ANTES de que se dibuje nada.
%  · positioning        → sintaxis «below left=2pt and 3pt».
%  · decorations.pathreplacing → llaves ({brace}) para señalar rangos.
\usetikzlibrary{babel,positioning,decorations.pathreplacing}
\usepackage{graphicx}
% Circuitos: el trabajo de electrónica se hace hoy con micro:bit y sus sensores
% integrados (MakeCode, sin cableado), así que no se necesitan esquemáticos.
% Si algún día entran montajes con protoboard, basta descomentar esta línea:
% los diagramas .tex/.tikz declarados en `recursos:` ya se incrustan solos.
% \usepackage{circuitikz}
\usepackage[most]{tcolorbox}
\usepackage{tabularx}
\usepackage{array}
\usepackage{fancyhdr}
\usepackage{titlesec}
\usepackage[document]{ragged2e}  % bandera a la derecha en todo el cuerpo
\usepackage{enumitem}
\usepackage{microtype}
\usepackage{etoolbox}
\usepackage{needspace}
% hyperref al final del preambulo: metadatos del PDF (titulo, autor, idioma,
% familia), marcadores por estacion (\pdfbookmark) y enlaces sin recuadro.
\usepackage[hidelinks,
  pdftitle={El que mira: por qué veinte testigos ayudan menos que uno},
  pdfauthor={PhD. Álvaro Cárdenas Orozco},
  pdfsubject={Territorio interior · Educación socioemocional · Grado 8° · Momento 3 · MILC v3},
  pdflang={es-CO},
  bookmarksopen=true,bookmarksnumbered=false]{hyperref}

% Helvetica Neue no trae la flecha U+2192, asi que cada «→» escrito literalmente
% en el YAML se compilaba a NADA, en silencio (462 flechas en 61 guias: rutas del
% tipo «paleta → tipografia → lockup» salian rotas). Mapeamos el caracter a su
% version matematica, que si tiene glifo. Asi el autor puede seguir escribiendo
% «→» comodamente en el YAML.
\usepackage{newunicodechar}
\newunicodechar{→}{\ensuremath{\rightarrow}}
% Mismo problema con los simbolos de marca y vinetas: el «✓» de «una palomita ✓»
% salia como cuadrito vacio en el PDF de todas las guias que lo usaban.
\usepackage{pifont}
\newunicodechar{✓}{\ding{51}}
\newunicodechar{✗}{\ding{55}}
\newunicodechar{✔}{\ding{52}}
\newunicodechar{•}{\textbullet}
\newunicodechar{≥}{\ensuremath{\geq}}
\newunicodechar{≤}{\ensuremath{\leq}}

\setmainfont{Helvetica Neue}
\setsansfont{Helvetica Neue}
\setlength{\parindent}{0pt}
\setlength{\parskip}{6pt}
% Sin particion de palabras: en 6.º-7.º una palabra cortada por guion al final
% de linea («Comuni-cacion») frena la lectura mucho mas que una linea corta.
\hyphenpenalty=10000
\exhyphenpenalty=10000
\pretolerance=2000
\tolerance=3000
\setlength{\emergencystretch}{3em}
\linespread{1.10}
\renewcommand{\arraystretch}{1.25}
\setlist{leftmargin=5mm,itemsep=1mm,topsep=1mm}
\newcolumntype{Y}{>{\RaggedRight\arraybackslash}X}

\definecolor{milcMagenta}{HTML}{E6007E}
\definecolor{milcVino}{HTML}{5A0038}
\definecolor{milcTurquesa}{HTML}{0093A5}
\definecolor{milcVerde}{HTML}{58A923}
\definecolor{milcMostaza}{HTML}{E5B400}
\definecolor{milcOcre}{HTML}{B86600}
\definecolor{milcNegro}{HTML}{241816}
\definecolor{milcGris}{HTML}{F7F7F4}
\definecolor{milcLinea}{HTML}{E8E2DD}
\definecolor{milcGradePrimary}{HTML}{0D9488}
\definecolor{milcGradeDark}{HTML}{115E59}
\definecolor{milcGradeSoft}{HTML}{CCFBF1}
\definecolor{milcGradeAccent}{HTML}{CFE7FF}
\definecolor{milcGradeText}{HTML}{FFFFFF}
\color{milcNegro}

\pagestyle{fancy}
\fancyhf{}
% Cabecera de dos lineas: identidad de la guia arriba y, a la derecha y en
% gris, la estacion en curso (\markright desde \milcestacion). Antes de la
% primera estacion la segunda linea va vacia; el \strut mantiene la altura
% para que la linea de arriba no baile entre paginas.
\lhead{\small Territorio interior · Educación socioemocional\\[-1pt]{\footnotesize\strut}}
\rhead{\small Grado 8° · Momento 3 · MILC v3\\[-1pt]{\footnotesize\strut\color{milcNegro!65}\rightmark}}
\cfoot{\small\thepage}
\renewcommand{\headrulewidth}{0pt}

% Sobre mostaza (#E5B400) y verde (#58A923) el texto va en milcNegro; sobre el
% resto de la paleta, en blanco. Se usa dentro de fonttitle/fontupper, donde un
% \color posterior manda sobre coltitle.
\newcommand{\milctintasobre}[1]{%
  \ifboolexpr{test{\ifstrequal{#1}{milcMostaza}} or test{\ifstrequal{#1}{milcVerde}}}%
    {\color{milcNegro}}{\color{white}}}

\newtcolorbox{titlebox}[1]{
  enhanced,colback=#1,colframe=#1,arc=7mm,boxrule=0pt,
  left=7mm,right=7mm,top=3mm,bottom=3mm,
  before skip=8pt,after skip=8pt,
  fontupper=\sffamily\bfseries\Large\color{white}
}

\newtcolorbox{softbox}[3]{
  enhanced,breakable,pad at break=3mm,
  colback=#2,colframe=#1,borderline west={4pt}{0pt}{#1},
  arc=2mm,boxrule=.4pt,left=4mm,right=4mm,top=3mm,bottom=3mm,
  before skip=6pt,after skip=8pt,title={#3},coltitle=white,
  colbacktitle=#1,fonttitle=\sffamily\bfseries\milctintasobre{#1},fontupper=\RaggedRight,
  attach boxed title to top left={xshift=3mm,yshift=-2mm},
  boxed title style={arc=2mm, boxrule=0pt}
}

\newtcolorbox{drawbox}[2]{
  enhanced,breakable,pad at break=4mm,
  colback=white,colframe=#1,arc=4mm,boxrule=1pt,
  left=5mm,right=5mm,top=4mm,bottom=4mm,before skip=8pt,after skip=8pt,
  title={#2},coltitle=white,colbacktitle=#1,fonttitle=\sffamily\bfseries\milctintasobre{#1},
  fontupper=\RaggedRight,
  attach boxed title to top left={xshift=4mm,yshift=-2mm},
  boxed title style={arc=3mm, boxrule=0pt}
}

\newtcolorbox{infoband}[1]{
  enhanced,breakable,pad at break=2mm,colback=#1!8,colframe=#1!50,arc=2mm,boxrule=.5pt,
  left=4mm,right=4mm,top=2mm,bottom=2mm,before skip=4pt,after skip=4pt,
  fontupper=\small\RaggedRight
}

% Puente narrativo entre secciones (contrato editorial Conéctate).
% Da continuidad: "Acabas de X. Ahora vas a Y porque Z."
% Visualmente: banda izquierda en color del grado, texto en itálica pequeña.
\newtcolorbox{puentebox}{
  enhanced,breakable,pad at break=2mm,colback=milcGradeSoft,colframe=milcGradeDark,
  borderline west={3pt}{0pt}{milcGradeDark},
  arc=0mm,boxrule=0pt,left=5mm,right=4mm,top=2.5mm,bottom=2.5mm,
  before skip=4pt,after skip=8pt,
  fontupper=\itshape\small\color{milcNegro}\RaggedRight
}
% La flecha va en modo matematico a proposito: Helvetica Neue NO tiene el glifo
% U+2192, asi que \textrightarrow se compilaba a NADA («Missing character: There
% is no → in font Helvetica Neue») y el puente salia sin flecha. En modo
% matematico la flecha viene de la fuente matematica y si se dibuja.
\newcommand{\puente}[1]{%
  \begin{puentebox}%
    \textcolor{milcGradeDark}{$\rightarrow$}\hspace{2mm}#1%
  \end{puentebox}%
}

\newcommand{\linead}[1]{\noindent\color{milcLinea}\rule{#1}{0.5pt}\color{milcNegro}}
\newcommand{\respuesta}[1]{\par\vspace{2mm}\foreach \n in {1,...,#1}{\noindent\color{milcLinea}\rule{\linewidth}{0.5pt}\par\vspace{5mm}}\color{milcNegro}}
\newcommand{\checkbox}{\textcolor{milcGradeDark}{\fbox{\phantom{x}}}\hspace{5pt}}

% ─── Listas aireadas para las guías socioemocionales ────────────────────
% Componer las verificaciones como «\checkbox texto\\» deja el texto que salta
% de línea debajo del cuadrito y apelmaza el bloque. Estos entornos alinean la
% continuación y airean los ítems. Los usa Territorio Interior; cualquier
% familia puede hacerlo.
\newlist{checklista}{itemize}{1}
\setlist[checklista]{
  label=\textcolor{milcGradeDark}{\fbox{\phantom{x}}},
  leftmargin=9mm, labelsep=3.2mm, itemsep=2.4mm,
  topsep=2.5mm, parsep=0pt, partopsep=0pt, before=\RaggedRight
}
\newlist{pasoslista}{enumerate}{1}
\setlist[pasoslista]{
  label=\textcolor{milcGradeDark}{\sffamily\bfseries\arabic*.},
  leftmargin=8mm, labelsep=3mm, itemsep=2.8mm,
  topsep=1mm, parsep=0pt, partopsep=0pt, before=\RaggedRight
}

\newcommand{\phasecard}[4]{
\begin{tcolorbox}[enhanced,colback=#3,colframe=#2,arc=3mm,boxrule=.5pt,height=3.05cm,valign=center,halign=center,left=1.5mm,right=1.5mm,top=2mm,bottom=2mm]
{\sffamily\bfseries\fontsize{22}{24}\selectfont\textcolor{#2}{#1}}\par
{\sffamily\bfseries\small #4}
\end{tcolorbox}
}

% ── Piezas del rediseno 2026 (legibilidad 6.º-11.º) ───────────────────
% Banda de estacion: reemplaza el titlebox macizo. Titulo a la izquierda,
% dato practico (tiempo/modalidad) a la derecha, en una sola linea.
% Nunca huerfana: pide 9 lineas libres (o salta de pagina rellenando con
% \vfill) y prohibe el salto justo despues de la banda. Nueve y no seis:
% la caja partible que sigue necesita ~80pt (titulo adosado, relleno y dos
% lineas) para arrancar en la misma pagina; con menos, tcolorbox la manda
% entera a la siguiente y la banda se queda sola. El penalty 10000 va
% en `after`, ANTES del espacio posterior: un `after skip` (aunque sea 0pt)
% es un \vskip tras la caja, y ese glue es un punto de corte legal que un
% \nopagebreak escrito despues de \end{tcolorbox} ya no cierra. Cada banda
% deja marcador PDF y actualiza la cabecera derecha.
\newcounter{milcestacion}
\newcommand{\milcestacion}[2]{%
\Needspace*{9\baselineskip}%
\stepcounter{milcestacion}%
\pdfbookmark[1]{#2}{milcestacion\themilcestacion}%
\markright{#2}%
\begin{tcolorbox}[enhanced,colback=#1,colframe=#1,arc=2mm,boxrule=0pt,
  left=5mm,right=5mm,top=2.6mm,bottom=2.6mm,before skip=11pt,
  after={\par\nopagebreak\vskip7pt}]
{\sffamily\bfseries\fontsize{15}{18}\selectfont\milctintasobre{#1}#2}
\end{tcolorbox}}

% Tarjeta de paso numerada: sustituye las tablas «Pilar / Aplicacion», que
% eran formato de consulta docente, no de estudiante. El numero va en un
% circulo sobre el borde para que la secuencia se lea de un vistazo.
\newcommand{\milcpaso}[3]{%
\Needspace{4\baselineskip}%
\begin{tcolorbox}[enhanced,colback=white,colframe=milcLinea,arc=1.5mm,boxrule=.9pt,
  left=14mm,right=4mm,top=3mm,bottom=3mm,before skip=4pt,after skip=4pt,
  fontupper=\RaggedRight,
  overlay={\node[anchor=north west,circle,fill=#2,text=white,inner sep=0pt,
    minimum size=7.4mm,font=\sffamily\bfseries\small]
    at ([xshift=3.6mm,yshift=-3.2mm]frame.north west) {#1};}]
#3
\end{tcolorbox}}

% Casilla marcable de tamano util con lapiz (la anterior era \fbox{\phantom{x}},
% de alto variable segun el texto que la seguia).
\newcommand{\milcmarca}{\framebox[3.8mm]{\rule{0pt}{3.4mm}}}

% Fila marcable de lista suelta (checks de la estacion 1).
\newcommand{\milccheckrow}[1]{%
\par\vspace{1.8mm}\noindent
\makebox[6.4mm][l]{\raisebox{-.55mm}{\textcolor{milcNegro}{\milcmarca}}}%
\parbox[t]{\dimexpr\linewidth-6.4mm\relax}{\RaggedRight #1}\par}

% Celda marcable dentro de la rubrica (tabla de autoevaluacion). Misma
% sangria francesa, pero pensada para vivir dentro de una columna X.
\newcommand{\milcopcion}[1]{%
\makebox[5.2mm][l]{\raisebox{-.55mm}{\textcolor{milcNegro}{\milcmarca}}}%
\parbox[t]{\dimexpr\linewidth-5.2mm\relax}{\RaggedRight #1}}

% ── Recursos visuales de la guía ──────────────────────────────────────
% Declarados en el bloque `recursos:` del YAML y emitidos por el builder.
% \guiaFigura   → imagen rasterizada o PDF (foto, carta celeste, montaje).
% \guiaDiagrama → fuente TikZ/circuitikz (.tex) incrustada como vector:
%                 es la vía para circuitos y esquemáticos editables.
% [#1] = texto alternativo (alt) · #2 = ruta relativa al .tex · #3 = pie
% El alt llega al PDF etiquetado (/Alt) y lo lee el lector de pantalla.
\newcommand{\guiaFigura}[3][]{%
\begin{center}
\includegraphics[width=0.88\linewidth,height=0.38\textheight,keepaspectratio,alt={#1}]{#2}\\[1.5mm]
{\footnotesize\itshape\textcolor{milcNegro!75}{#3}}
\end{center}
}

\newcommand{\guiaDiagrama}[2]{%
\begin{center}
\input{#1}\\[1.5mm]
{\footnotesize\itshape\textcolor{milcNegro!75}{#2}}
\end{center}
}

\begin{document}
\thispagestyle{empty}

% ── Portada ───────────────────────────────────────────────────────────
% Antes: un rectangulo de color a pagina completa. Se veia bien en pantalla,
% pero en la fotocopiadora del colegio se comia el toner y salia gris sucio.
% Ahora la banda ocupa el tercio superior y el resto va en blanco.
\begin{tikzpicture}[remember picture,overlay]
  \path[fill=milcGradePrimary]
    (current page.north west) rectangle ([yshift=-10.6cm]current page.north east);

  \node[anchor=north west,text=milcGradeText] at ([xshift=2.0cm,yshift=-1.55cm]current page.north west)
    {\sffamily\bfseries\fontsize{12}{14}\selectfont MOMENTO};
  \node[anchor=north west,text=milcGradeText,opacity=.85] at ([xshift=2.0cm,yshift=-2.35cm]current page.north west)
    {\sffamily\bfseries\fontsize{8.5}{10}\selectfont TERRITORIO INTERIOR · SOCIOEMOCIONAL};
  \node[anchor=north east,text=milcGradeText,opacity=.85,align=right] at ([xshift=-2.0cm,yshift=-1.55cm]current page.north east)
    {\sffamily\bfseries\fontsize{9}{12}\selectfont I.E. Sor María Juliana\\Cartago, Valle del Cauca};

  \node[anchor=west,text=milcGradeText] (gradonum) at ([xshift=1.85cm,yshift=-7.1cm]current page.north west)
    {\sffamily\bfseries\fontsize{112}{112}\selectfont 3};
  % El circulito de grado (°) solo aplica a las guias de grado («11°»); en
  % Bebras y Semillero el numero grande es un momento/modulo, no un grado.
  % Se posiciona relativo al nodo `gradonum`, no en coordenadas absolutas.
  

  \node[anchor=west,text=milcGradeText,text width=11.4cm,align=left] at ([xshift=1.5cm,yshift=0.5cm]gradonum.east)
    {
     {\sffamily\bfseries\fontsize{9}{11}\selectfont\MakeUppercase{Territorio interior · Socioemocional}}\\[3mm]
     {\sffamily\bfseries\fontsize{21}{25}\selectfont\setlength{\baselineskip}{27pt}El que\\[4pt]mira\\[4pt]veinte ayudan menos que uno}\\[3.5mm]
     {\sffamily\bfseries\fontsize{15}{18}\selectfont Grado 8° · Momento 3}
    };
\end{tikzpicture}

\vspace*{9.4cm}
\vfill

{\sffamily\bfseries\fontsize{8}{10}\selectfont\textcolor{milcGradeDark}{LO QUE TE LLEVAS HOY}}

\begin{tcolorbox}[enhanced,colback=white,colframe=milcNegro,arc=2mm,boxrule=1.6pt,
  left=5mm,right=5mm,top=4mm,bottom=4mm,before skip=3pt,after skip=12pt,
  fontupper=\RaggedRight\fontsize{11}{15.5}\selectfont]
Tu plan de espectador: las tres cosas que puedes hacer sin ser el héroe —nombrar a alguien, distraer y acompañar después—, escritas y ensayadas, más la frase exacta con que romperías la difusión de responsabilidad.
\end{tcolorbox}

\begin{tcolorbox}[enhanced,colback=milcGris,colframe=milcGris,arc=2mm,boxrule=0pt,
  left=5mm,right=5mm,top=3.5mm,bottom=3.5mm,after skip=12pt]
{\sffamily\bfseries\fontsize{8}{10}\selectfont\textcolor{milcNegro!55}{ESTA GUÍA ES DE}}\\[3mm]
\begin{tabularx}{\linewidth}{X p{3.2cm} p{3.2cm}}
\linead{\linewidth} & \linead{3.0cm} & \linead{3.0cm}\\[-1mm]
{\fontsize{8.5}{10}\selectfont\textcolor{milcNegro!55}{Nombre completo}} &
{\fontsize{8.5}{10}\selectfont\textcolor{milcNegro!55}{Grupo}} &
{\fontsize{8.5}{10}\selectfont\textcolor{milcNegro!55}{Fecha}}\\
\end{tabularx}
\end{tcolorbox}

\vfill

\noindent\textcolor{milcLinea}{\rule{\linewidth}{1pt}}\\[2mm]
{\sffamily\bfseries\fontsize{9.5}{12}\selectfont Autor: Dr. Álvaro Cárdenas Orozco}\\[1mm]
{\sffamily\fontsize{8.5}{11}\selectfont\textcolor{milcNegro!55}{Metodología MILC · Apertura ancestral · El papel del espectador · Triángulo Dussel-Estoicismo-Floridi}}

\newpage

\section*{Apertura: El que mira: por qué veinte testigos ayudan menos que uno}

\begin{softbox}{milcGradeDark}{milcGradeSoft}{Saber ancestral}
En Colombia hay una figura que casi nadie conoce y que sirve exactamente para esto: \textbf{el juez de paz}. La Constitución de 1991 la creó en su \textbf{artículo 247}, y la \textbf{Ley 497 de 1999} la reglamentó. \textbf{Fíjate en cómo funciona,} porque es lo contrario de lo que uno esperaría de un juez. \textbf{No hay que ser abogado} ---ni tener ningún título en derecho---. \textbf{Lo eligen los vecinos por votación popular}, entre gente de la propia comunidad, y hay que \textbf{haber vivido allí al menos un año}. \textbf{Y no falla según los códigos: falla \emph{en equidad}}, es decir, según los criterios de justicia de esa comunidad. \emph{Lo que se le pide es una sola cosa: ser alguien en quien la gente confía.} \textbf{Al lado existe otra figura parecida, el \emph{conciliador en equidad}}, también un vecino, también sin toga. \textbf{Piensa en lo que eso dice.} \emph{En un conflicto del barrio, el que tiene autoridad para meterse no es el que tiene un cargo: es el que la comunidad señaló como confiable.} \textbf{La autoridad, ahí, no viene de arriba: viene de al lado.}
\end{softbox}

\begin{softbox}{milcTurquesa}{milcGris}{Saber contemporáneo}
¿Y por qué, si veinte personas ven algo injusto, no interviene ninguna? \textbf{Se midió, y el resultado es de los más incómodos de la psicología.} En 1968, \textbf{John Darley} y \textbf{Bibb Latané} pusieron a estudiantes a conversar por un intercomunicador. En medio de la conversación, uno de los participantes ---grabado--- \textbf{sufría lo que parecía un ataque}. Cada persona creía que había \textbf{una, dos o cinco} más escuchando. \textbf{Los números son brutales.} Cuando la persona creía ser \textbf{la única} que oía, \textbf{el 85\,\% avisó}. Cuando creía que había \textbf{otra} más, bajó al \textbf{62\,\%}. Cuando creía que había \textbf{cinco}, \textbf{solo el 31\,\% avisó}. \textbf{Y no es que los demás fueran peores personas: también tardaban más.} El tiempo de reacción subió de \textbf{52 segundos} a \textbf{93} y luego a \textbf{166} (Darley y Latané, 1968). \textbf{El nombre del fenómeno lo explica todo: \emph{difusión de responsabilidad}}. \emph{Cuando hay mucha gente, cada uno siente que le toca una fracción del deber ---y una fracción no alcanza para levantarse---.} \textbf{Todos piensan «alguien más va a hacer algo». Y como todos lo piensan, no lo hace nadie.}
\end{softbox}

\begin{softbox}{milcMostaza}{milcGris}{Pregunta-puente}
\textit{El juez de paz tiene autoridad \textbf{porque los vecinos se la dieron, no porque tenga un cargo}. \textbf{¿Cuántas veces este año esperaste a que «alguien» hiciera algo} ---un profesor, el monitor, el amigo del afectado--- \textbf{en vez de ser tú ese alguien}?}
\end{softbox}

\begin{softbox}{milcVerde}{milcGris}{Saber y saber hacer}
Al terminar podrás: (1) \textbf{explicar} la difusión de responsabilidad con los datos del experimento, y reconocerla en tu propio salón; (2) \textbf{entender} que la autoridad para intervenir no viene de un cargo, como muestran el juez de paz y el conciliador en equidad; (3) \textbf{crear} tres formas de intervenir que no exigen enfrentarse a nadie ---nombrar, distraer y acompañar después---; (4) \textbf{usar} la frase que rompe la difusión: señalar a \textbf{una} persona concreta en vez de pedirle ayuda al aire.
\end{softbox}



\small
\begin{tabularx}{\textwidth}{>{\bfseries}p{3.0cm}Y}
\rowcolor{milcGradePrimary!12}Autor & Dr. Álvaro Cárdenas Orozco\\
DBA & Comprende los fenómenos que afectan a su territorio, regula sus emociones ante ellos y acompaña a otros con empatía (transversal 6°--11°, articulado con la política «Escuela, territorio de vida» del MEN, Resolución 006519 de 2025, y la Ley 1620 de 2013)\\
Referentes & Primera ayuda psicológica (OMS, 2011/2012) · Directrices IASC (2007) · USGS y Servicio Geológico Colombiano · Pueblo nasa y la minga (CRIC) · Dussel · Estoicismo · Floridi\\
Producto final & Tu plan de espectador: las tres cosas que puedes hacer sin ser el héroe —nombrar a alguien, distraer y acompañar después—, escritas y ensayadas, más la frase exacta con que romperías la difusión de responsabilidad.\\
\end{tabularx}
\normalsize

\puente{En el momento 2 conseguiste un aliado para \textbf{tus} decisiones. Hoy le damos la vuelta: \textbf{qué haces cuando el que está en problemas es otro} y tú solo estás mirando.}

\milcestacion{milcGradeDark}{Ruta MILC: así vamos a aprender}
\begin{tabularx}{\textwidth}{>{\centering\arraybackslash}X>{\centering\arraybackslash}X>{\centering\arraybackslash}X>{\centering\arraybackslash}X}
\phasecard{1}{milcVerde}{milcGris}{Escucha\\[-1mm]\small Observo, escucho y pregunto.} &
\phasecard{2}{milcTurquesa}{milcGris}{Sistematización\\[-1mm]\small Miro y decido.} &
\phasecard{3}{milcMagenta}{milcGris}{Praxis\\[-1mm]\small Construyo y aplico.} &
\phasecard{4}{milcVino}{milcGris}{Evaluación\\[-1mm]\small Reflexión liberadora.}
\end{tabularx}


\puente{Primero, tus propias veces de público: cuando viste algo y no hiciste nada. \emph{Le pasa a todo el mundo, y el experimento explica por qué.}}

\milcestacion{milcVerde}{Estación 1 de 3 · Escucha}

\begin{softbox}{milcVerde}{milcGris}{Escena para escuchar}
\textbf{Actividad 1 · IDENTIFICA --- Las veces que fui público.} \textbf{Nadie va a leer esta página.} \textbf{Primera parte.} Recuerda \textbf{dos situaciones de este año} en que viste algo que estaba mal ---una burla que se pasó, alguien excluido, algo que se reenvió--- \textbf{y no hiciste nada}. Escribe qué pasó y \textbf{cuánta gente había mirando}. \textbf{Segunda parte, la clave.} ¿\textbf{Qué te dijiste} en el momento? \emph{«No es conmigo» · «alguien va a decir algo» · «si me meto me cae a mí» · «de pronto entendí mal» · «ya casi se acaba».} \textbf{Escríbelo textual}, porque esas frases son el mecanismo. \textbf{Tercera parte.} ¿Alguna vez alguien \textbf{sí} hizo algo? ¿Qué hizo exactamente, y qué pasó después? \textbf{Y la última:} ¿alguna vez \textbf{tú} fuiste el que estaba en problemas y nadie hizo nada? \emph{Anota qué habrías querido que hiciera alguien ---no lo heroico: lo mínimo---.}
\end{softbox}

{\sffamily\bfseries\fontsize{8}{10}\selectfont\textcolor{milcNegro!55}{MARCA CUANDO LO HAYAS HECHO}}
\milccheckrow{Escribiste dos situaciones con \textbf{cuánta gente había mirando}.}
\milccheckrow{Anotaste textual la frase con que te lo justificaste en el momento.}
\milccheckrow{Escribiste qué habrías querido que alguien hiciera por ti, en lo mínimo.}

\begin{infoband}{milcVerde}


\textbf{En tu cuaderno (Actividad 1 · IDENTIFICA).} Título: «Actividad 1. Las veces que fui público». \textbf{Formato:} las dos situaciones con el número de testigos + la frase textual + lo que habrías querido que hicieran por ti. \textbf{Extensión:} media página. \textbf{Sabes que terminaste cuando} tienes escrita \textbf{la frase exacta} con que te quedaste quieto. Esta página es tuya: \textbf{nadie más tiene que leerla si tú no quieres}.
\end{infoband}

\puente{Ya las tienes. Ahora los números ---que son peores de lo que uno cree--- y por qué la autoridad para meterse no es un cargo.}

\milcestacion{milcTurquesa}{Fase 2 · Sistematización: entender al que mira}

\begin{softbox}{milcTurquesa}{milcGris}{Cuatro claves sobre mirar}
Cuatro claves para entender por qué el público no ayuda ---y qué lo convierte otra vez en gente.
\end{softbox}

\milcpaso{1}{milcTurquesa}{\textbf{Entre más somos, menos ayudamos. Con números.} En el experimento de Darley y Latané, quien creía ser \textbf{el único} que escuchaba avisó el \textbf{85\,\%} de las veces; con \textbf{una} persona más, el \textbf{62\,\%}; con \textbf{cinco} más, apenas el \textbf{31\,\%} (Darley y Latané, 1968). \emph{Y el tiempo también se estira:} de 52 segundos a 166. \textbf{Léelo al derecho:} \emph{tú tienes más probabilidad de recibir ayuda si hay una sola persona mirando que si hay veinte}. \textbf{Es contra toda intuición y es la clave de todo lo demás:} el problema del salón lleno \textbf{no es que a nadie le importe}; es que \textbf{el deber se repartió en pedacitos} y ningún pedacito alcanza para moverse.}
\milcpaso{2}{milcTurquesa}{\textbf{Nadie es «el que debería», y por eso hay que nombrarlo.} La difusión de responsabilidad se rompe de una manera casi tonta: \textbf{señalando a una persona}. \emph{«¡Que alguien llame al coordinador!» no mueve a nadie ---le habla al aire---.} \textbf{«Tú, el de la camisa azul: llama al coordinador» sí}, porque le devuelve a una persona el cien por ciento del deber. \emph{Esto sirve en los dos papeles:} si eres el que mira, \textbf{nombra}; si eres el que está en problemas, \textbf{no pidas ayuda en general: pídesela a alguien con nombre propio}. \textbf{Y es exactamente lo que hace un aviso bien hecho:} convertir a un público en una persona.}
\milcpaso{3}{milcTurquesa}{\textbf{La autoridad para meterse no es un cargo.} La razón más común para quedarse quieto es \emph{«yo quién soy para decir algo»}. \textbf{El juez de paz existe para desmentir eso:} es un vecino, sin título de abogado, elegido por votación, que resuelve \textbf{en equidad} según los criterios de justicia de su comunidad ---y lo único que se le exige es \textbf{ser alguien confiable} (Ley 497 de 1999)---. \emph{La autoridad, ahí, no baja de un escritorio: la pone la gente de al lado.} \textbf{En un salón pasa lo mismo:} no hace falta ser el monitor ni el profesor para decir «bájenle». \textbf{Lo que da peso a esa frase no es un cargo: es que alguien la diga.}}
\milcpaso{4}{milcTurquesa}{\textbf{Tres cosas que puedes hacer, y ninguna es enfrentarse.} La razón por la que casi nadie interviene es que \textbf{solo se imagina una forma de hacerlo}: pararse de frente y encarar. \emph{Eso da miedo, y con razón.} \textbf{Pero hay tres más, y sirven igual o más.} \textbf{Nombrar:} señalar a alguien concreto para que actúe, o avisarle a un adulto. \textbf{Distraer:} interrumpir la escena sin mencionarla ---«oye, ¿me prestas el cuaderno?», «vámonos que ya sonó»---; \emph{le quita el público al que está haciendo el daño, y sin público casi nada sigue}. \textbf{Acompañar después:} buscar a la persona cuando ya pasó y decirle dos frases. \emph{Esta última es la más fácil de todas y la que más recuerda el que la recibe.}}

\begin{drawbox}{milcTurquesa}{Las tres intervenciones, de menor a mayor riesgo}
\begin{pasoslista}
\item \textbf{Acompañar después.} Dos frases, en privado, cuando ya pasó. \emph{Riesgo casi cero, y es lo que más se recuerda.}
\item \textbf{Distraer.} Interrumpir sin nombrar el asunto: pedir algo, cambiar de tema, irse llevándose a alguien. \emph{Le quita el público.}
\item \textbf{Nombrar.} «Tú, \_\_\_\_: ven un momento» · «profe, ¿puede venir?». \emph{Rompe la difusión: le devuelve a una persona todo el deber.}
\item \textbf{Y solo si es seguro:} decirlo de frente, corto y sin sermón ---«ya, con eso no»---.
\item \textbf{Si hay riesgo físico, no se negocia:} es de las situaciones del momento 5, y ahí se avisa a un adulto.
\end{pasoslista}
\end{drawbox}

\begin{softbox}{milcMostaza}{milcGris}{Errores comunes y su corrección}
\textbf{«No es conmigo»:} el público es parte de la escena; sin público casi nada continúa. \textbf{Pedir ayuda al aire:} «que alguien haga algo» no le habla a nadie. \textbf{Creer que solo vale enfrentarse:} es la forma más riesgosa y la menos usada; hay tres más. \textbf{Grabar en vez de intervenir:} el celular convierte al testigo en público. \textbf{Esperar al profesor:} a veces hay que llamarlo ---nombrándolo---, pero esperar es la difusión en acción. \textbf{«Yo quién soy»:} el juez de paz tampoco es abogado. \textbf{Intervenir con sermón:} pone al grupo del lado del que hizo el daño.
\end{softbox}

\begin{infoband}{milcTurquesa}


\textbf{En tu cuaderno (Actividad 2 · EXPLICA).} Título: «Actividad 2. La difusión de responsabilidad». \textbf{Formato:} escribe los tres porcentajes del experimento y explica con tus palabras por qué bajan; después vuelve a una de tus dos situaciones y escribe \textbf{cuál de las tres intervenciones habría cabido ahí}. \textbf{Extensión:} media página. \textbf{Sabes que terminaste cuando} puedes explicar por qué \textbf{es mejor tener un solo testigo que veinte}.
\end{infoband}

\puente{Ya sabes qué pasa. Es hora de tener tres acciones listas, ninguna de las cuales exige ser valiente ni enfrentarse a nadie.}

\milcestacion{milcMagenta}{Fase 3 · Praxis: dejar de ser público}

\begin{softbox}{milcMagenta}{milcGris}{Cómo se interviene sin ser el héroe}
Ahora se arman las tres acciones. La idea no es volverte valiente: es que tengas algo listo para los cuatro segundos en que hay que decidir.
\end{softbox}

\milcpaso{1}{milcMagenta}{\textbf{Nombrar (8 min).} Escribe \textbf{tu frase para señalar a alguien}, en dos versiones: una para pedirle a un compañero que actúe y otra para avisarle a un adulto. \textbf{Condición:} que \textbf{contenga un nombre} ---o algo que identifique sin duda a una persona---. \emph{Pruébalas leyéndolas en voz alta y preguntándote: si me lo dijeran a mí, ¿sabría que me habla a mí?}}
\milcpaso{2}{milcMagenta}{\textbf{Distraer (8 min).} Escribe \textbf{tres maneras de interrumpir una escena sin mencionarla}: pedir algo prestado, llamar a alguien, recordar una hora, proponer irse. \emph{Tienen que ser cosas normales, que no delaten que estás interviniendo ---esa es justamente la gracia---.} \textbf{Es la herramienta más subestimada:} sin público, casi nada sigue.}
\milcpaso{3}{milcMagenta}{\textbf{Acompañar después (8 min).} Escribe las \textbf{dos frases} que le dirías en privado a alguien a quien le pasó algo. \textbf{Sin consejo, sin «deberías», sin preguntar por qué no se defendió.} \emph{Recuerda lo del momento 1 del grado 7: escuchar, no arreglar.} \textbf{Y una regla:} no le prometas que vas a hacer algo si no lo vas a hacer.}
\milcpaso{4}{milcMagenta}{\textbf{El ensayo de los segundos (12 min, en grupos de cuatro).} Uno describe una escena ---\textbf{inventada, nunca real y nunca sobre alguien presente}---, dos hacen de público y uno interviene. \textbf{Lo que se cronometra es una sola cosa: cuánto tarda en hacer algo.} Roten hasta que todos pasen. \emph{Después comparen: ¿qué intervención salió más fácil? Casi siempre gana \textbf{distraer}, y casi nadie la había pensado antes.}}

\begin{drawbox}{milcMagenta}{Antes de cerrar tu plan}
\begin{checklista}
\item Tu frase para \textbf{nombrar} contiene un nombre o algo que identifique sin duda a alguien.
\item Tienes \textbf{tres maneras de distraer} que parecen cosas normales.
\item Tus dos frases para \textbf{acompañar después} no traen consejo ni reproche.
\item Sabes cuál de las tres usar según el riesgo, y que enfrentarse es la última.
\item Entendiste que si hay riesgo físico eso \textbf{no es asunto de espectador}: se avisa a un adulto.
\item Ensayaste en voz alta y mediste \textbf{cuánto tardas} en hacer algo.
\end{checklista}
\end{drawbox}

\begin{softbox}{milcMagenta}{milcGris}{Plantilla de trabajo}
\textbf{Nombrar.} \emph{«\_\_\_\_ (nombre), ven un momento» · «Profe \_\_\_\_, ¿puede venir?»} ---nunca «que alguien\ldots»---. \\[1mm]
\textbf{Distraer.} \emph{«¿Me prestas \_\_\_\_?» · «Vámonos que ya sonó» · «\_\_\_\_ te está buscando».} \\[1mm]
\textbf{Acompañar después.} \emph{«Vi lo que pasó. Me pareció una mala jugada.»} ---dos frases, sin consejo---.
\end{softbox}

\begin{infoband}{milcMagenta}


\textbf{En tu cuaderno (Actividad 3 · CREA).} Título: «Actividad 3. Mi plan de espectador». \textbf{Formato:} la frase de nombrar (dos versiones) + las tres maneras de distraer + las dos frases de acompañar + cuánto tardaste en el ensayo. \textbf{Extensión:} una página. \textbf{Sabes que terminaste cuando} tienes \textbf{tres acciones distintas} y sabes cuál usarías según el riesgo.
\end{infoband}

\puente{El producto es un plan de espectador. \emph{No convierte a nadie en héroe: convierte al público en gente que hace algo pequeño.}}

\milcestacion{milcMostaza}{Producto de la sesión}
\textbf{Plan de espectador: nombrar, distraer y acompañar después, ensayados y ordenados por riesgo}.
\begin{softbox}{milcMostaza}{milcGris}{Criterios de calidad}
(1) La frase para nombrar identifica sin ambigüedad a una persona, no le habla al aire. (2) Las tres maneras de distraer son verosímiles y no delatan la intervención. (3) Las frases para acompañar después no dan consejo ni reprochan a quien recibió el daño. (4) El plan ordena las intervenciones por riesgo y reserva el enfrentamiento como última opción. (5) Se reconoce el límite: ante riesgo físico se avisa a un adulto, no se resuelve entre pares.
\end{softbox}

\puente{Antes de cerrar, revisa lo importante: si tu único plan es «meterme de frente», no lo vas a usar nunca, y por eso hay tres.}

\milcestacion{milcVino}{Evaluación liberadora · cinco dimensiones}

\begin{softbox}{milcVino}{milcGris}{Sentido de la evaluación}
Evaluar no es solo revisar una nota. Es reconocer qué aprendí, qué cambió en mí, qué emociones manejé, cómo me relaciono con la ciudad y la comunidad, y qué le dejaré a quien viene detrás.
\end{softbox}

\textbf{Mi reflexión liberadora en cinco dimensiones.} Completa cada frase iniciada:

\small
\begin{tabularx}{\textwidth}{>{\bfseries\RaggedRight\arraybackslash}p{3.4cm}Y>{\RaggedRight\arraybackslash}p{4.6cm}}
\rowcolor{milcVino}\color{white}Dimensión & \color{white}\bfseries Mi respuesta (frame starter) & \color{white}\bfseries Algo que descubrí\\
1. Personal & Lo que más cambió en mí fue \linead{6.5cm} & \linead{4.2cm}\\
2. Emocional & La emoción más fuerte fue \linead{2.5cm} y la regulé \linead{3cm} & \linead{4.2cm}\\
3. Ciudadana & En democracia esto importa porque \linead{6cm} & \linead{4.2cm}\\
4. Local (Cartago) & En mi barrio sería útil para \linead{6.5cm} & \linead{4.2cm}\\
5. Intergeneracional & A mi abuela/abuelo le diría \linead{6.5cm} & \linead{4.2cm}\\
\end{tabularx}
\normalsize

\begin{infoband}{milcVino}
\textbf{Para la próxima sustentación.} Vuelve a esta tabla en seis meses. Verás que las respuestas cambian — no porque lo aprendido cambie, sino porque tú cambias. La evaluación liberadora es una conversación contigo a través del tiempo.
\end{infoband}


\puente{Cerramos con tres voces: el rostro que reclama un derecho, la disposición a que te muestren que estabas mal, y qué se cultiva en un lugar compartido.}

\milcestacion{milcGradeDark}{Triángulo de pensamiento · cierre filosófico}

\begin{softbox}{milcVino}{milcGris}{Dussel · lente del nosotros}
\textit{«El otro se revela realmente como otro\ldots como el pobre, el oprimido; el que, a la vera del camino, fuera del sistema, muestra su rostro sufriente y sin embargo desafiante: «¡Tengo hambre!, ¡tengo derecho a comer!»»}

{\footnotesize\itshape\color{milcNegro!70}--- Enrique Dussel · Filosofía de la liberación (1977)}

Dussel dice que el otro \textbf{se revela}: aparece, se hace ver. \emph{Y el problema del espectador es exactamente ese ---que ya lo vio---.} \textbf{Después de mirar, quedarse quieto ya no es neutral:} es una respuesta. \textbf{Y fíjate otra vez en el final de la cita:} el que sufre \textbf{reclama un derecho}, no pide un favor. \emph{Eso cambia la pregunta que uno se hace:} no es «¿me corresponde a mí ser bueno con esta persona?», \textbf{es «¿estoy presenciando que le quiten algo que le corresponde?»}. \emph{A la primera pregunta uno puede contestar que no le nace. A la segunda, no.}

\textbf{Lo último que vi y dejé pasar, ¿era «algo feo» o era alguien perdiendo un derecho?}
\end{softbox}
\linead{\linewidth}\\[1mm]
\linead{\linewidth}

\begin{softbox}{milcMostaza}{milcGris}{Estoicismo · Marco Aurelio · Meditaciones VI, 21 (c. 175 d.C.) · cuidado interior}
\textit{«Si alguien puede refutarme y probar de modo concluyente que pienso o actúo incorrectamente, de buen grado cambiaré de proceder. Pues persigo la verdad, que no dañó nunca a nadie; en cambio, sí se daña el que persiste en su propio engaño e ignorancia.»}

{\footnotesize\itshape\color{milcNegro!70}--- Marco Aurelio · Meditaciones VI, 21 (c. 175 d.C.)}

Aquí hay dos lecciones y las dos sirven hoy. \textbf{La primera es para el que interviene:} conviene hacerlo \emph{como quien corrige}, no como quien condena ---corto, sin sermón---, porque el que se siente humillado se atrinchera. \textbf{La segunda es para uno mismo:} \emph{«se daña el que persiste en su propio engaño»}. \textbf{Y el engaño del espectador tiene un nombre concreto: «no es conmigo».} \emph{Repetírselo veinte veces no lo vuelve verdad ---lo vuelve costumbre---,} y la costumbre es lo que hace que a los dos años uno ya ni note lo que está pasando al lado.

\textbf{¿Cuántas veces he dicho «no es conmigo» este mes? ¿Sigo notando lo que pasa?}
\end{softbox}
\linead{\linewidth}\\[1mm]
\linead{\linewidth}

\begin{softbox}{milcTurquesa}{milcGris}{Floridi · lente de la infoesfera}
\textit{«El florecimiento de las entidades informacionales y de toda la infoesfera debe promoverse preservando, cultivando y enriqueciendo sus propiedades.»}

{\footnotesize\itshape\color{milcNegro!70}--- Luciano Floridi · La ética de la información (2013; trad.)}

Hoy el espectador tiene una opción que antes no existía y que \textbf{empeora las cosas: grabar}. \emph{Grabar se siente como hacer algo ---uno está «documentando»--- y en realidad es la forma más pura de ser público:} añade audiencia, y encima vuelve el momento permanente. \textbf{La regla es simple:} \emph{si vas a sacar el celular, que sea para llamar a alguien, no para filmar}. \textbf{Y del lado de cultivar:} el equivalente digital de «acompañar después» es el mensaje privado, \textbf{y en un chat cuesta todavía menos} porque nadie te ve mandarlo.

\textbf{La última vez que alguien grabó algo en mi salón, ¿para qué sirvió esa grabación?}
\end{softbox}
\linead{\linewidth}\\[1mm]
\linead{\linewidth}

\begin{infoband}{milcNegro}
\textbf{Nota para el docente.} Las tres son citas textuales y llevan su referencia debajo. Puedes remitir a la obra en clase.
\end{infoband}

\milcestacion{milcVino}{Mi autoevaluación · marca una casilla por fila}

{\small
\setlength{\extrarowheight}{2.2pt}
\begin{tabularx}{\textwidth}{>{\RaggedRight\arraybackslash\bfseries}p{3.0cm}YYY}
\rowcolor{milcVino}\color{white}\bfseries Criterio & \color{white}\bfseries Lo logré & \color{white}\bfseries En proceso & \color{white}\bfseries Necesito apoyo\\[.6mm]
Comprensión del tema & \milcopcion{Explico las ideas con mis palabras.} & \milcopcion{Reconozco las ideas, falta precisión.} & \milcopcion{Necesito repasar lo básico.}\\[.8mm]
\rowcolor{milcGris}Calidad de la intervención & \milcopcion{Tengo tres acciones distintas y sé cuál usar según el riesgo.} & \milcopcion{Sé que debería hacer algo, pero solo se me ocurre meterme de frente.} & \milcopcion{Sigo pensando que si no soy yo el que hace el daño, no es asunto mío.}\\[.8mm]
Aplicación y producto & \milcopcion{Mi producto cumple los criterios.} & \milcopcion{Está armado, con ajustes pendientes.} & \milcopcion{Aún está en construcción.}\\[.8mm]
\rowcolor{milcGris}Reflexión 5 dimensiones & \milcopcion{Respondí cada una con honestidad.} & \milcopcion{Respondí algunas con detalle.} & \milcopcion{Mis respuestas son muy generales.}\\[.8mm]
Triángulo de pensamiento & \milcopcion{Conecté las 3 voces con mi trabajo.} & \milcopcion{Conecté al menos una.} & \milcopcion{No las relaciono con mi proceso.}\\[.8mm]
\end{tabularx}}

\vspace{3mm}
\textbf{Hoy entendí que:}\\[1mm]
\linead{\linewidth}\\[1mm]
\linead{\linewidth}

\textbf{Mi compromiso para la próxima sesión es:}\\[1mm]
\linead{\linewidth}\\[1mm]
\linead{\linewidth}

\begin{softbox}{milcMostaza}{milcGris}{Cita de despedida · Marco Aurelio}
\textit{``El obstáculo en el camino se vuelve el camino. Lo que estorba al camino se vuelve camino.''}\\[1mm]
Cada error de hoy, bien observado, se convierte en pieza del camino que pisarás mañana.
\end{softbox}

\small
\begin{tabularx}{\textwidth}{>{\bfseries}p{3.6cm}Y}
\rowcolor{milcGradeDark!12}Nombre completo: & \linead{10cm}\\
Fecha: & \linead{4cm} \quad Curso/grupo: \linead{4cm}\\
Mi firma: & \linead{9cm}\\
\end{tabularx}
\normalsize

\begin{infoband}{milcGradeDark}
\textbf{Para la próxima sesión.} Dos cosas. \textbf{Primero, usa una de las tres} ---la más fácil sirve: \textbf{acompañar después}---. Busca a alguien a quien le haya pasado algo esta semana y dile tus dos frases. \emph{Anota qué te respondió; te va a sorprender lo mucho que pesa algo tan pequeño.} \textbf{Segundo, pregunta en tu casa quién arreglaba los líos:} averigua con alguien mayor \textbf{quién intervenía cuando había un problema entre vecinos} ---si había un juez de paz, un conciliador, el presidente de la Junta de Acción Comunal, el párroco, una señora respetada de la cuadra---, \textbf{cómo llegó esa persona a ese papel} y si la gente le hacía caso. \textbf{Trae:} qué te respondieron y lo que te contaron. \emph{Vas a encontrar lo mismo casi siempre: nadie la nombró desde arriba. Era alguien de al lado en quien la gente confiaba.}
\end{infoband}

\end{document}
